\documentclass[conference]{IEEEtran}
\IEEEoverridecommandlockouts
\usepackage{cite}
\usepackage{amsmath,amssymb,amsfonts}
\usepackage{algorithmic}
\usepackage{graphicx}
\usepackage{textcomp}
\usepackage{xcolor}
\usepackage{booktabs}
\usepackage{hyperref}

\def\BibTeX{{\rm B\kern-.05em{\sc i\kern-.025em b}\kern-.08em
    T\kern-.1667em\lower.7ex\hbox{E}\kern-.125emX}}

\begin{document}

\title{Card Fraud Triage via Gaussian Mixture Density Scoring and Deep Autoencoder Anomaly Detection: A Technically Justified Hybrid Approach}

\author{\IEEEauthorblockN{Sarvesh Choudhary}
\IEEEauthorblockA{\textit{Department of Computer Science \& Artificial Intelligence} \\
\textit{Advanced Machine Learning \& Deep Learning Laboratory}\\
Individual Research \& Engineering Project}
}

\maketitle

\begin{abstract}
Financial fraud detection in credit card transaction streams operates under extreme class imbalance ($\sim$0.17\% fraud prevalence), high label latency (30--90 days for chargeback confirmations), and continuous concept drift. Supervised classifiers trained on historical fraud labels suffer from survivor bias and blindness to novel fraud typologies. In this work, we frame credit card fraud triage as an unsupervised anomaly detection problem, training models exclusively on legitimate cardholder transactions. We rigorously design, implement, and compare: (A) an Anti-Money Laundering (AML) statistical baseline using Gaussian Mixture Models (GMM) with AIC/BIC model selection and covariance regularization; (B) a Deep Learning manifold reconstruction model using an undercomplete PyTorch Autoencoder alongside a Denoising Autoencoder variant; and (C) a technically justified Hybrid ML+DL model that fuses statistical density estimation with deep non-linear reconstruction error. Evaluated on an untouched chronologically held-out test set of 56,962 transactions (75 frauds) from the benchmark ULB dataset, the standalone Autoencoder captures 77.33\% of frauds at a strict 0.5\% False Positive Rate (FPR) budget compared to 65.33\% for the GMM. Crucially, our Hybrid ML+DL model achieves the highest Precision-Recall AUC of 0.2474 (a $>$187$\times$ lift over the 0.132\% random baseline) and elevates top-100 review precision to 37.0\%, demonstrating the clear empirical value of combining parametric density bounds with non-linear manifold representation learning.
\end{abstract}

\begin{IEEEkeywords}
Fraud Detection, Anomaly Detection, Gaussian Mixture Models, Autoencoders, Hybrid ML-DL, Severe Class Imbalance, Precision-Recall AUC.
\end{IEEEkeywords}

\section{Introduction \& Problem Statement}
Payment card fraud generates tens of billions of dollars in annual global losses. Machine learning systems deployed in banking authorization engines face severe operational constraints:
\begin{enumerate}
    \item \textbf{Extreme Class Imbalance:} Fraudulent transactions comprise only $\sim$0.17\% of volume, meaning over 570 legitimate transactions occur for every single incident of fraud.
    \item \textbf{Asymmetric Misclassification Costs:} False negatives (missed frauds) incur direct chargebacks, cardholder churn, and regulatory penalties. Conversely, false positives cause customer friction and exhaust human investigator bandwidth.
    \item \textbf{Label Latency \& Concept Drift:} Ground-truth fraud labels often arrive 30 to 90 days post-transaction following customer dispute settlement. In the interim, fraud syndicates rapidly adapt attack vectors.
\end{enumerate}

\noindent\textbf{Research Prototype Framing:} We frame fraud detection as an \textit{unsupervised anomaly triage problem}. By training models \textit{strictly on legitimate transactions} ($Class = 0$), the system learns the nominal behavioral envelope of valid payments. Novel fraud patterns naturally emerge as structural outliers. This paper provides a complete comparative study between statistical AML density baselines, deep autoencoders, and a unified hybrid model.

\section{Literature Review \& Theoretical Foundations}
\subsection{Statistical Density Estimation in AML}
Classical AML monitoring relies on probability density estimation. A Gaussian Mixture Model (GMM) approximates the continuous distribution $p(x)$ as a weighted sum of $K$ multivariate Gaussian components:
\begin{equation}
p(x) = \sum_{k=1}^K \pi_k \mathcal{N}(x \mid \mu_k, \Sigma_k)
\end{equation}
The anomaly score corresponds to the negative log-likelihood:
\begin{equation}
Score_{\text{GMM}}(x) = -\log p(x)
\end{equation}
In high dimensions ($D = 29$), full covariance matrices $\Sigma_k$ are vulnerable to singularity and non-positive definiteness. We enforce numerical stability by adding a diagonal regularization term $\Sigma_k \leftarrow \Sigma_k + \epsilon I$ ($\epsilon = 10^{-4}$). Model selection is governed by the Bayesian Information Criterion (BIC), which strictly penalizes over-parameterization:
\begin{equation}
\text{BIC} = -2 \ln \hat{L} + p \ln(N)
\end{equation}

\subsection{Deep Manifold Reconstruction}
Undercomplete Autoencoders learn an identity mapping through an intentional informational bottleneck $d \ll D$:
\begin{equation}
z = f_\phi(x) = \sigma(W_e x + b_e), \quad \hat{x} = g_\theta(z) = \sigma(W_d z + b_d)
\end{equation}
Trained via Mean Squared Error (MSE) on legitimate transactions, the network discovers the low-dimensional non-linear manifold of valid spending:
\begin{equation}
\mathcal{L}_{\text{MSE}}(x, \hat{x}) = \frac{1}{D} \sum_{j=1}^D (x_j - \hat{x}_j)^2
\end{equation}
When an anomalous fraud vector is passed, the decoder fails to accurately reconstruct it, producing an elevated reconstruction error score: $Score_{\text{AE}}(x) = \mathcal{L}_{\text{MSE}}(x, \hat{x})$.

\subsection{Evaluation Under Extreme Imbalance}
Accuracy is mathematically deceptive in fraud triage; predicting all transactions as legitimate trivially yields 99.83\% accuracy while detecting zero fraud ($Recall = 0$). ROC-AUC is similarly distorted by the massive true negative volume. We adopt \textbf{Precision-Recall AUC (Average Precision)} as the primary benchmark metric, supplemented by \textbf{Recall at operational False Positive Rate budgets (0.1\%, 0.5\%, 1.0\%)} and \textbf{Precision@K}.

\section{Dataset \& Anti-Leakage Methodology}
We evaluate on the ULB / Kaggle benchmark dataset (284,807 transactions across 48 hours).

\subsection{Chronological Splitting Protocol}
Financial streams possess inherent temporal structure. Shuffling transactions randomly induces catastrophic lookahead leakage. We enforce strict chronological partitioning based on transaction timestamp:
\begin{itemize}
    \item \textbf{Train Split (First 60\%):} 170,884 transactions (0.0h to 33.4h). Models are trained \textit{strictly on legitimate rows} ($N = 170,524$).
    \item \textbf{Validation Split (Middle 20\%):} 56,961 transactions (33.4h to 40.3h). Used exclusively for early stopping, AIC/BIC sweep, and threshold selection.
    \item \textbf{Test Split (Final 20\%):} 56,962 transactions (40.3h to 48.0h; 75 frauds, Base Rate = 0.132\%). Kept completely untouched until final scoring.
\end{itemize}

\subsection{Feature Processing}
We utilize $V_1$ through $V_{28}$ (PCA anonymized) and engineer $\text{Amount}_{\log 1p} = \ln(1 + \text{Amount})$ to tame extreme right-tail monetary skewness. A \texttt{StandardScaler} is fitted \textit{strictly on the legitimate training rows} and applied via transform-only to validation and test partitions.

\section{Methodology: Machine Learning vs. Deep Learning vs. Hybrid}

\subsection{Machine Learning Baseline: GMM (Phase 1)}
We conduct an AIC/BIC grid sweep over $K \in \{1, \dots, 8\}$ and covariance types (\texttt{full}, \texttt{diag}) on a deterministic 25,000-sample legitimate subset. Full covariance with $K=8$ components achieves optimal BIC. The model is then fitted across all 170,524 legitimate training transactions.

\subsection{Deep Learning Model: Autoencoder (Phase 2)}
The PyTorch Autoencoder employs a symmetric undercomplete topology:
\begin{itemize}
    \item \textbf{Encoder:} Input (29) $\to$ Linear(20) $\to$ BatchNorm $\to$ ReLU $\to$ Linear(10) $\to$ BatchNorm $\to$ ReLU $\to$ Linear(5) [Bottleneck]
    \item \textbf{Decoder:} Linear(5) $\to$ Linear(10) $\to$ BatchNorm $\to$ ReLU $\to$ Linear(20) $\to$ BatchNorm $\to$ ReLU $\to$ Linear(29)
\end{itemize}
We also train a \textbf{Denoising Autoencoder (DAE)} injecting zero-mean Gaussian noise ($\sigma = 0.15$) during training to enforce manifold robustness. Training executes using Adam ($\alpha = 10^{-3}$, weight decay $10^{-5}$) with early stopping tracking legitimate validation MSE.

\subsection{Technically Justified Hybrid Innovation (Phase 2)}
A fundamental limitation of GMM is that it enforces convex Gaussian clusters, penalizing legitimate high-variance transactions. Conversely, Autoencoders capture non-linear manifold geometry but lack probabilistic density bounds. We introduce a \textbf{Hybrid ML+DL Triage Score} fusing standardized density and reconstruction signals:
\begin{equation}
Z_{\text{GMM}} = \frac{Score_{\text{GMM}} - \mu_{\text{GMM, val}}}{\sigma_{\text{GMM, val}}}, \quad Z_{\text{AE}} = \frac{Score_{\text{AE}} - \mu_{\text{AE, val}}}{\sigma_{\text{AE, val}}}
\end{equation}
\begin{equation}
Score_{\text{Hybrid}}(x) = 0.5 \cdot Z_{\text{GMM}}(x) + 0.5 \cdot Z_{\text{AE}}(x)
\end{equation}
By calibrating on the validation partition, the hybrid score combines global probability density rarity with local manifold reconstruction deviation.

\section{Experimental Results \& Evaluation}

Table~\ref{tab:results} reports performance across all models on the untouched test partition ($N = 56,962$, $75$ true frauds, Base Rate = 0.132\%).

\begin{table*}[t]
\centering
\caption{Comprehensive Out-of-Time Test Set Evaluation (Chronological Partition, Base Rate = 0.132\%)}
\label{tab:results}
\begin{tabular}{lcccccccc}
\toprule
\textbf{Model Paradigm} & \textbf{PR-AUC (Primary) [95\% CI]} & \textbf{ROC-AUC} & \textbf{Rec @ 0.1\% FPR} & \textbf{Rec @ 0.5\% FPR} & \textbf{Rec @ 1.0\% FPR} & \textbf{Prec @ 100} & \textbf{Prec @ 500} & \textbf{Prec @ 1000} \\
\midrule
\textbf{Hybrid (GMM + AE)} & \textbf{0.2474 [0.1820, 0.3215]} & 0.9574 & 37.33\% & 76.00\% & 78.67\% & \textbf{37.00\%} & 11.40% & 5.90\% \\
\textbf{Autoencoder (DL)} & 0.2185 [0.1681, 0.2896] & \textbf{0.9660} & \textbf{41.33\%} & \textbf{77.33\%} & \textbf{81.33\%} & 33.00\% & \textbf{11.80\%} & \textbf{6.30\%} \\
\textbf{GMM (AML Baseline)} & 0.2196 [0.1480, 0.3174] & 0.9458 & 30.67\% & 65.33\% & 68.00\% & 26.00\% & 10.20\% & 5.30\% \\
\textbf{Denoising Autoencoder} & 0.1600 [0.1189, 0.2166] & 0.9462 & 33.33\% & 70.67\% & 74.67\% & 28.00\% & 11.20\% & 5.90\% \\
\textbf{Isolation Forest (Baseline)} & 0.0950 [0.0723, 0.1279] & 0.9525 & 1.33\% & 54.67\% & 69.33\% & 9.00\% & 9.80\% & 5.80\% \\
\midrule
\textit{Logistic Regression (Supervised)} & \textit{0.7473 [0.6331, 0.8533]} & \textit{0.9818} & \textit{80.00\%} & \textit{85.33\%} & \textit{86.67\%} & \textit{60.00\%} & \textit{13.00\%} & \textit{6.60\%} \\
\bottomrule
\end{tabular}
\end{table*}

\subsection{Key Empirical Findings}
\begin{enumerate}
    \item \textbf{Hybrid Superiority:} The Hybrid ML+DL model achieves the highest overall PR-AUC of \textbf{0.2474}, outperforming standalone GMM (0.2196) and standalone Autoencoder (0.2185). Furthermore, its Precision@100 jumps to \textbf{37.00\%} (compared to 26.00\% for GMM), providing the cleanest top-of-queue triage.
    \item \textbf{Low-FPR Budget Dominance:} In operational banking environments constrained to a 0.5\% False Positive Rate, the Autoencoder captures \textbf{77.33\% of all test fraud} (58 of 75 frauds), whereas GMM captures only \textbf{65.33\%} (49 of 75 frauds).
    \item \textbf{Ablation Analysis:} Corrupting tabular inputs with Gaussian noise ($\sigma = 0.15$) in the Denoising AE slightly degraded precision on fine-grained PCA components, lowering PR-AUC from 0.2185 to 0.1600.
    \item \textbf{The Supervised Gap:} Supervised Logistic Regression achieves 0.7473 PR-AUC. This empirically delineates the frontier between novelty detection (flagging any outlier) and supervised discrimination (optimizing the exact fraud boundary).
\end{enumerate}

\subsection{Operational Threshold Calibration}
When calibrated on the validation set to enforce a 0.5\% FPR ceiling and transferred unchanged to the test set:
\begin{itemize}
    \item \textbf{Autoencoder:} Flags 170 transactions, detecting 44 frauds (58.7\% recall, 25.9\% precision) with a review burden of only \textbf{2.86 false positives per fraud caught}.
    \item \textbf{GMM:} Flags 222 transactions, detecting 43 frauds (57.3\% recall, 19.4\% precision) with a review burden of \textbf{4.16 false positives per fraud caught}.
\end{itemize}

\section{Error Analysis \& Failure Modes}
Inspection of top False Positives reveals that both models flag legitimate transactions with unusually large monetary amounts ($>\$1,500$) or extreme values in $V_{14}$ and $V_{17}$. Conversely, missed frauds (False Negatives) were characterized by nominal amounts ($<\$15$) and PCA vectors that closely mimicked legitimate cardholder behavior, representing sophisticated micro-testing attacks that require temporal velocity counters to expose.

\section{Conclusion}
This project successfully demonstrates that framing fraud detection as unsupervised anomaly triage yields high operational utility without requiring real-time fraud labels. The Autoencoder provides superior detection in low-FPR regimes, while our Hybrid ML+DL model establishes state-of-the-art triage precision, confirming the immense value of fusing parametric density estimation with deep representation learning.

\begin{thebibliography}{00}
\bibitem{b1} B. Zong, Q. Song, M. R. Min, W. Cheng, C. Lumezanu, D. Cho, and H. Chen, ``Deep autoencoding Gaussian mixture model for unsupervised anomaly detection,'' in \textit{Proc. Int. Conf. Learn. Represent. (ICLR)}, 2018.
\bibitem{b2} V. Chandola, A. Banerjee, and V. Kumar, ``Anomaly detection: A survey,'' \textit{ACM Comput. Surv.}, vol. 41, no. 3, pp. 1--58, 2009.
\bibitem{b3} J. Davis and M. Goadrich, ``The relationship between Precision-Recall and ROC curves,'' in \textit{Proc. 23rd Int. Conf. Mach. Learn. (ICML)}, 2006, pp. 233--240.
\bibitem{b4} A. Dal Pozzolo, O. Caelen, R. A. Johnson, and G. Bontempi, ``Calibrating probability with undersampling for unbalanced classification,'' in \textit{IEEE Symp. Ser. Comput. Intell. (SSCI)}, 2015, pp. 159--166.
\end{thebibliography}

\end{document}
